% !TEX program = xelatex
\documentclass[11pt,a4paper]{article}
\usepackage{algo-preamble}

\title{\textbf{L04 \textperiodcentered{} Διαίρει \& Κυρίευε II}\\[2pt]
\large Inversions, Dominant Colour \& Karatsuba}
\author{Algorithms Class Hub}
\date{\small Αλγόριθμοι και Πολυπλοκότητα (K17) \textperiodcentered{} ΕΚΠΑ}

\begin{document}
\maketitle

\begin{center}\itshape\small
Τρία προβλήματα που λύνονται κομψά με διαίρει και κυρίευε: μέτρηση αντιστροφών σε
$O(n \log n)$, εύρεση κυρίαρχου χρώματος σε $O(n^2 \log n)$, και ο πολλαπλασιασμός
Karatsuba σε $O(n^{1.585})$.
\end{center}

\section{Τι θα δούμε}

Στο προηγούμενο μάθημα είδαμε το \textbf{σχήμα} D\&C και το χρησιμοποιήσαμε για
ταξινόμηση. Εδώ θα δείξουμε ότι το ίδιο σχήμα --- διαίρει, κυρίευσε αναδρομικά,
συνδύασε --- λύνει τρία εντελώς διαφορετικά προβλήματα:

\begin{enumerate}
  \item \textbf{Μέτρηση αντιστροφών} σε πίνακα --- από $O(n^2)$ πέφτουμε σε
  $O(n \log n)$.
  \item \textbf{Κυρίαρχο χρώμα} σε σκακιέρα --- από $O(n^4)$ πέφτουμε σε
  $O(n^2 \log n)$.
  \item \textbf{Πολλαπλασιασμός μεγάλων αριθμών} --- από $O(n^2)$ πέφτουμε σε
  $O(n^{1.585})$ με τον αλγόριθμο \textbf{Karatsuba}.
\end{enumerate}

Σε κάθε ένα από τα τρία θα ξαναδείς το ίδιο μοτίβο: μια \textbf{έξυπνη σύνθεση}
των απαντήσεων των υποπροβλημάτων που κερδίζει χρόνο. Το Master Theorem από το
προηγούμενο μάθημα θα είναι το κύριο εργαλείο ανάλυσης.

\begin{keypoint}
\textbf{Το μάθημα της διάλεξης:} «Διαίρει και κυρίευε» δεν είναι μια εξίσωση,
είναι ένας \textbf{τρόπος σκέψης}. Όταν δεις πρόβλημα στο εξεταζόμενο μάθημα που
μοιάζει άγνωστο, η πρώτη ερώτηση είναι: «μπορώ να το σπάσω σε μικρότερα ίδια
προβλήματα και να συγχωνεύσω τις λύσεις σε γραμμικό χρόνο;» Αν ναι --- Master
Theorem.
\end{keypoint}

\section{Σύντομη υπενθύμιση: Master Theorem}

Από το προηγούμενο μάθημα, αν
\[
T(n) = a \, T\!\left(\tfrac{n}{b}\right) + O(n^d), \qquad a > 0,\ b \ge 1,\ d
\ge 0
\]
τότε:
\[
T(n) = \begin{cases}
O(n^d)          & \text{αν } d > \log_b a \\
O(n^d \log n)   & \text{αν } d = \log_b a \\
O(n^{\log_b a}) & \text{αν } d < \log_b a
\end{cases}
\]

\textbf{Διαισθητικά:} βρίσκεις τη βάση $b$ (πόσες φορές μικρότερο γίνεται το
υποπρόβλημα), τον αριθμό $a$ (πόσα τέτοια υποπροβλήματα έχεις) και το $d$ (κόστος
συγχώνευσης). Συγκρίνεις $a$ με $b^d$ και κερδίζει είτε η ρίζα (μεγάλο combine),
είτε όλο το δέντρο εξίσου (ίσιο), είτε τα φύλλα (πολλά υποπροβλήματα).

\begin{intuition}
\textbf{Σήμα κινδύνου:} το Master Theorem ισχύει μόνο όταν η αναδρομή έχει
σταθερό λόγο διαχωρισμού $n/b$. Αναδρομές όπως $T(n) = T(n/3) + T(2n/3) + n$ ή
$T(n) = T(n-1) + n$ δεν λύνονται έτσι --- επιστρέφουμε στις τεχνικές του
προηγούμενου μαθήματος (δέντρο αναδρομής, telescoping, ισχυρή επαγωγή).
\end{intuition}

\subsection{Γενικότερη μορφή (Cormen)}

Στη γενικότερη διατύπωση το $O(n^d)$ αντικαθίσταται από οποιαδήποτε συνάρτηση
$f(n)$:
\begin{itemize}
  \item Αν $f(n) = O(n^{\log_b a - \epsilon})$ για κάποια $\epsilon > 0$:
  $T(n) = \Theta(n^{\log_b a})$.
  \item Αν $f(n) = \Theta(n^{\log_b a})$: $T(n) = \Theta(n^{\log_b a} \log n)$.
  \item Αν $f(n) = \Omega(n^{\log_b a + \epsilon})$ και $a f(n/b) \le c f(n)$ για
  κάποιο $c < 1$ και αρκετά μεγάλο $n$: $T(n) = \Theta(f(n))$.
\end{itemize}
Στο εξεταζόμενο μάθημα η απλή μορφή είναι σχεδόν πάντα αρκετή.

\subsection{Ασκήσεις για να συντονιστούμε}

Δοκίμασε \textbf{νοερά} να βρεις τη λύση:
\begin{itemize}
  \item $T(n) = T(n-1) + 2$. (Όχι Master --- telescope: $T(n) = \Theta(n)$.)
  \item $T(n) = 2T(n-1) + 1$. (Όχι Master --- βγαίνει $\Theta(2^n)$, όπως ο
  πύργος του Hanoi.)
  \item $T(n) = 3T(n/4) + n$. ($a = 3,\ b = 4,\ d = 1$. $\log_4 3 \approx 0.79 <
  1$, άρα \textbf{case 1}: $T(n) = O(n)$.)
  \item $T(n) = T(n/3) + T(2n/3) + cn$. (\textbf{Όχι Master} --- άνισος
  διαχωρισμός. Δέντρο αναδρομής $\to \Theta(n \log n)$.)
  \item $T(n) = T(\sqrt{n}) + \log n$. (Αλλαγή μεταβλητής $n = 2^m$:
  $T(2^m) = T(2^{m/2}) + m$, που βγαίνει $\Theta(\log n \cdot \log \log n)$.)
\end{itemize}

\section{Πρόβλημα 1: Μέτρηση Αντιστροφών}

\subsection{Ορισμός}

\textbf{Είσοδος:} πίνακας $A$ με $n$ φυσικούς αριθμούς.

\textbf{Ζητούμενο:} το πλήθος των ζευγαριών $(i, j)$ με $i < j$ και
$A[i] > A[j]$.

\textbf{Παράδειγμα.} Στον πίνακα $[1, 3, 5, 2, 4]$ οι αντιστροφές είναι:
$(3, 2)$ --- το $3$ είναι πριν από το $2$ αλλά μεγαλύτερο --- καθώς και $(5, 2)$
και $(5, 4)$. Σύνολο: \textbf{3 αντιστροφές}.

\begin{intuition}
Διαισθητικά, οι αντιστροφές μετρούν \textbf{πόσο μακριά είναι ο πίνακας από το να
είναι ταξινομημένος}. Αν είναι ήδη ταξινομημένος $\to$ 0 αντιστροφές. Αν είναι
ταξινομημένος ανάποδα $\to \binom{n}{2}$ αντιστροφές (το μέγιστο).
\end{intuition}

\subsection{Γιατί μας νοιάζει}

Πρακτικές εφαρμογές που χρησιμοποιούν τη μέτρηση αντιστροφών ως μέτρο
\textbf{ομοιότητας κατάταξης}: συνεργατικό φιλτράρισμα σε εφαρμογές
μουσικής/ταινιών (πόσο μοιάζουν δύο κατατάξεις --- μετράμε τις αντιστροφές μεταξύ
τους), θεωρία ψηφοφορίας (απόσταση Kendall tau), μέτρηση «ποσοστού ταξινόμησης»
ενός πίνακα, και ευαισθησία της σειράς κατάταξης στις μηχανές αναζήτησης.

\begin{center}
\begin{tabular}{@{}l c c c c c@{}}
\toprule
     & Α & Β & Γ & Δ & Ε \\
\midrule
Εγώ  & 1 & 2 & 3 & 4 & 5 \\
Εσύ  & 1 & 3 & 5 & 2 & 4 \\
\bottomrule
\end{tabular}
\end{center}

\subsection{Αφελής λύση: $\Theta(n^2)$}

Διπλό loop:
\begin{codebox}
\begin{verbatim}
inversions(A):
    count <- 0
    Για i = 1 ως n-1:
        Για j = i+1 ως n:
            Εάν A[i] > A[j]: count <- count + 1
    επίστρεψε count
\end{verbatim}
\end{codebox}
Πολυπλοκότητα $\Theta(n^2)$. Για $n = 10^6$, $10^{12}$ συγκρίσεις $\to$ πρακτικά
αδύνατο.

\subsection{Λύση με D\&C: $O(n \log n)$}

\textbf{Σπάμε} τον πίνακα στη μέση και ορίζουμε \textbf{τρία είδη} αντιστροφών:
\begin{enumerate}
  \item Αντιστροφές που είναι \textbf{εντελώς μέσα} στο αριστερό μισό $A_1$.
  \item Αντιστροφές που είναι \textbf{εντελώς μέσα} στο δεξί μισό $A_2$.
  \item Αντιστροφές \textbf{«μικτές»}: το μεγαλύτερο στοιχείο στο $A_1$, το
  μικρότερο στο $A_2$.
\end{enumerate}
Τα δύο πρώτα είδη μετρώνται \textbf{αναδρομικά}. Το ζητούμενο είναι: πώς μετράμε
τα μικτά γρήγορα;

\begin{notebox}
Παράδειγμα στον πίνακα $[10, 5, 8, 6, 3, 4, 7, 1, 9, 2, 11]$: το αριστερό μισό
$A_1 = [10, 5, 8, 6, 3, 4]$ δίνει αναδρομικά $12$ αντιστροφές, το δεξί μισό
$A_2 = [7, 1, 9, 2, 11]$ δίνει $3$· οι μικτές αντιστροφές είναι $15$, άρα
συνολικά $12 + 3 + 15 = 30$.
\end{notebox}

\subsection{Πώς συνδυάζουμε σε γραμμικό χρόνο;}

Παρατήρηση-κλειδί: \textbf{αν οι δύο υποπίνακες είναι ταξινομημένοι} μπορούμε να
μετρήσουμε τις μικτές αντιστροφές \textbf{καθώς τους συγχωνεύουμε} σε έναν
μεγαλύτερο ταξινομημένο πίνακα. Αυτό μας θυμίζει τη \texttt{merge} της mergesort
από το προηγούμενο μάθημα --- και πραγματικά, ο αλγόριθμος \textbf{ταυτόχρονα}
ταξινομεί και μετράει.

\textbf{Ιδέα της \texttt{merge-and-count}:} σαρώνουμε τους δύο πίνακες με δύο
δείκτες $i$ (πάνω από $A_1$) και $j$ (πάνω από $A_2$). Κάθε φορά παίρνουμε το
μικρότερο. Αν παίρνουμε το στοιχείο από το $A_2$ (δηλαδή το «μπλε» είναι
μικρότερο από το «κόκκινο»), τότε \textbf{όλα τα υπόλοιπα στοιχεία του $A_1$ είναι
μεγαλύτερα από αυτό} --- άρα μεγαλύτερα \textbf{και} σε προγενέστερη θέση $\to$
όλα σχηματίζουν αντιστροφή με αυτό το στοιχείο. Προσθέτουμε στον μετρητή.

\begin{keypoint}
Επειδή το $A_1$ είναι \textbf{ταξινομημένο}, αν το μπλε στοιχείο είναι μικρότερο
από το τρέχον κόκκινο στη θέση $i$, αυτό σημαίνει ότι είναι αυτόματα μικρότερο
και από όλα τα κόκκινα στις θέσεις $i, i+1, \dots, m$. Αυτά είναι $m - i + 1$
στοιχεία. Γι' αυτό προσθέτουμε ακριβώς αυτή την ποσότητα σε μία πράξη --- όχι
έναν-έναν.
\end{keypoint}

\subsection{Ο αλγόριθμος \texttt{merge-and-count}}

\begin{codebox}
\begin{verbatim}
Διαδικασία merge-and-count(a, b):
  // a, b ταξινομημένοι πίνακες μεγέθους m και n
  r <- 0                      // πλήθος μικτών αντιστροφών
  i <- 1; j <- 1
  Για k = 1 ως m + n:
    Εάν a[i] < b[j]:
      c[k] <- a[i];  i <- i + 1
    αλλιώς:
      c[k] <- b[j];  j <- j + 1
      r <- r + (m - i + 1)    // όλα τα υπόλοιπα του a είναι > b[j]
  επίστρεψε (r, c)
\end{verbatim}
\end{codebox}
\textbf{Πολυπλοκότητα:} $O(m + n)$ --- μία γραμμική σάρωση. Πανομοιότυπη με την
\texttt{merge} της mergesort, με μόνη προσθήκη τη γραμμή
\texttt{r <- r + (m - i + 1)}.

\subsection{Ο αλγόριθμος \texttt{sort-and-count}}

\begin{codebox}
\begin{verbatim}
Διαδικασία sort-and-count(A):
  Εάν |A| = 1:
    επίστρεψε (0, A)
  Χώρισε τον A σε δύο μισά A1, A2
  (r1, A1) <- sort-and-count(A1)
  (r2, A2) <- sort-and-count(A2)
  (r,  A ) <- merge-and-count(A1, A2)
  επίστρεψε (r1 + r2 + r, A)
\end{verbatim}
\end{codebox}
Επιστρέφει \textbf{και} τον αριθμό αντιστροφών \textbf{και} τον ταξινομημένο
πίνακα --- χρειαζόμαστε το ταξινομημένο για να μπορεί η \texttt{merge-and-count}
του γονέα να δουλέψει.

\subsection{Πολυπλοκότητα}

\[
T(n) = 2 T(n/2) + O(n)
\]
Ακριβώς η ίδια αναδρομή με την mergesort. Master Theorem: $a = 2,\ b = 2,\
d = 1$, $\log_2 2 = 1 = d$ $\to$ \textbf{case 2}: $T(n) = O(n \log n)$.

\textbf{Έλεγχος στο παράδειγμα.} Στον πίνακα $[1, 3, 5, 2, 4]$: αριστερό
$[1, 3, 5] \to 0$ αντιστροφές, ταξινομημένο· δεξί $[2, 4] \to 0$ αντιστροφές,
ταξινομημένο. Συγχώνευση: παίρνω 1, παίρνω 2 (κερδίζει το δεξί ενώ απομένουν 2
στο αριστερό από τη θέση 2 $\to +2$), παίρνω 3, παίρνω 4 (κερδίζει ενώ απομένει 1
$\to +1$), παίρνω 5. \textbf{Σύνολο: 3.} $\checkmark$

\section{Πρόβλημα 2: Κυρίαρχο Χρώμα}

\subsection{Ορισμός}

\textbf{Είσοδος:} σκακιέρα $n \times n$ (όπου $n = 2^k$) με κάθε πλακίδιο
χρωματισμένο.

\textbf{Ζητούμενο:} να βρούμε αν υπάρχει χρώμα που εμφανίζεται σε
\textbf{περισσότερα από τα μισά} πλακίδια (δηλαδή σε $> n^2 / 2$ θέσεις). Αν ναι,
ποιο είναι.

\begin{notebox}
\textbf{Δεν} μας λένε πόσα χρώματα υπάρχουν συνολικά. Μπορεί να είναι 2, μπορεί
να είναι $n^2$. Αυτό κάνει το πρόβλημα μη τετριμμένο: δεν μπορούμε απλά να
φτιάξουμε πίνακα μετρητών μήκους «αριθμός χρωμάτων» χωρίς να ξέρουμε πόσοι είναι.
\end{notebox}

\textbf{Παραδείγματα:} σκακιέρα $4 \times 4$ με 9 κόκκινα πλακίδια από 16 $\to$
κυρίαρχο χρώμα $=$ κόκκινο ($9 > 8$). Σκακιέρα $4 \times 4$ με 8 κόκκινα $\to$
\textbf{όχι} κυρίαρχο ($8 \not> 8$ --- αυστηρή ανισότητα).

\subsection{Αφελής λύση: $O(n^4)$}

Σαρώνουμε τα $n^2$ πλακίδια. Για κάθε νέο χρώμα κρατάμε μετρητή και τον
ενημερώνουμε. Αλλά \textbf{σε κάθε επίσκεψη} δεν ξέρουμε αν το χρώμα είναι νέο,
οπότε ψάχνουμε στη λίστα των ήδη γνωστών χρωμάτων ($m$ χρώματα $\Rightarrow O(m)$
ανά πλακίδιο). Συνολικά $O(m \cdot n^2)$, και αφού $m \le n^2$:
\[
T(n) = O(n^4)
\]
Αργό. (Με hash table πέφτει σε $O(n^2)$ αμορτοποιημένα --- αλλά εδώ θέλουμε να
δούμε αν το D\&C βελτιώνει.)

\subsection{Λύση με D\&C: $O(n^2 \log n)$}

\textbf{Διαίρει:} σπάμε τη σκακιέρα σε 4 υποσκακιέρες $\frac{n}{2} \times
\frac{n}{2}$.

\textbf{Κυρίευσε:} καλούμε αναδρομικά. Κάθε υποσκακιέρα επιστρέφει είτε ένα
\textbf{υποψήφιο κυρίαρχο χρώμα} $c_i$, είτε «κανένα».

\textbf{Συνδύαζε:} με την εξής παρατήρηση-κλειδί:

\begin{keypoint}
\textbf{Αν η αρχική σκακιέρα έχει κυρίαρχο χρώμα $c$, τότε το $c$ είναι κυρίαρχο
σε τουλάχιστον μία από τις 4 υποσκακιέρες.}

\textbf{Απόδειξη:} Έστω ότι το $c$ είναι κυρίαρχο στην αρχική, άρα εμφανίζεται
$> n^2/2 = 4 \cdot (n/2)^2 / 2$ φορές. Αν το $c$ \textbf{δεν} ήταν κυρίαρχο σε
\textbf{καμία} υποσκακιέρα, τότε σε κάθε μία εμφανίζεται \textbf{το πολύ}
$(n/2)^2 / 2$ φορές. Άθροισμα: το πολύ $4 \cdot (n/2)^2 / 2 = n^2/2$ ---
αντίφαση. Άρα τουλάχιστον μία υποσκακιέρα έχει το $c$ ως κυρίαρχο.
$\blacksquare$
\end{keypoint}

\textbf{Αυτό κάνει τη ζωή μας εύκολη:} οι μόνοι υποψήφιοι κυρίαρχοι της αρχικής
σκακιέρας είναι τα (το πολύ) \textbf{4 χρώματα} που επιστρέφουν τα 4 αναδρομικά
καλέσματα. Για κάθε ένα από αυτά τα 4, μετράμε σε όλη τη σκακιέρα πόσες φορές
εμφανίζεται. Αυτό είναι $O(n^2)$ ανά χρώμα $\to O(4 n^2) = O(n^2)$ συνολικά.

\subsection{Πολυπλοκότητα}

\[
T(n) = 4 T(n/2) + O(n^2)
\]
Master Theorem: $a = 4,\ b = 2,\ d = 2$, $\log_2 4 = 2 = d$ $\to$ \textbf{case
2}: $T(n) = O(n^2 \log n)$.

\textbf{Σχόλιο: μπορεί καλύτερα;} Με hash table σε $O(n^2)$ μέσο όρο. Με
ντετερμινιστικό Boyer-Moore majority vote σε $O(n^2)$ ντετερμινιστικά (το
πρόβλημα ανάγεται στην 1D εκδοχή). Αλλά αυτό που μετράει εδώ είναι ότι \textbf{το
D\&C είναι μια καθαρή, καλά αναλυμένη λύση} και είναι το ακριβές πρότυπο που η
εξεταζόμενη μπορεί να ζητήσει --- διότι η αρχή της λύσης (παρατήρηση $+$ Master
Theorem) μετριέται σε όλη την ύλη.

\section{Πρόβλημα 3: Πολλαπλασιασμός Ακεραίων --- Karatsuba}

\subsection{Πόσο κοστίζει ο πολλαπλασιασμός;}

Δουλεύουμε με \textbf{πολύ μεγάλους αριθμούς} --- π.χ.\ αριθμοί $n$ ψηφίων όπου
$n = 10^4$ ή $10^9$. Δεν χωράνε σε \texttt{int64}. Πρέπει να τους
αναπαραστήσουμε ως πίνακες ψηφίων.

\begin{center}
\begin{tabular}{@{}l l@{}}
\toprule
Πράξη & Πολυπλοκότητα (σχολικός τρόπος) \\
\midrule
Πρόσθεση/πολλαπλασιασμός \textbf{μονοψήφιων} & $O(1)$ \\
Πρόσθεση \textbf{$n$-ψήφιων}                 & $O(n)$ \\
Πολλαπλασιασμός \textbf{$n$-ψήφιων}          & $O(n^2)$ \\
\bottomrule
\end{tabular}
\end{center}

Το $O(n^2)$ προέρχεται από τη σχολική μέθοδο: για κάθε ψηφίο του πρώτου αριθμού
πολλαπλασιάζουμε με όλα τα ψηφία του δεύτερου $\to n \times n$ μονοψήφιοι
πολλαπλασιασμοί.

\textbf{Παράδειγμα:} $1357 \cdot 2468$:
\[
\begin{array}{r}
1357 \\
\times\ 2468 \\
\hline
10856 \\
8142\phantom{0} \\
5428\phantom{00} \\
2714\phantom{000} \\
\hline
3{,}349{,}076
\end{array}
\]

\subsection{Πρώτη προσπάθεια D\&C --- αποτυγχάνει}

Σπάμε κάθε αριθμό στα μέσα ψηφία του: $1357 = 13 \cdot 10^2 + 57$ και
$2468 = 24 \cdot 10^2 + 68$. Γενικά:
\[
x = \alpha \cdot 10^{n/2} + \beta, \qquad y = \gamma \cdot 10^{n/2} + \delta
\]
όπου τα $\alpha, \beta, \gamma, \delta$ έχουν $n/2$ ψηφία. Τότε:
\[
x \cdot y = \alpha \gamma \cdot 10^n + (\alpha \delta + \beta \gamma) \cdot
10^{n/2} + \beta \delta
\]
Έχουμε \textbf{τέσσερις} πολλαπλασιασμούς μικρότερων αριθμών: $\alpha\gamma$,
$\alpha\delta$, $\beta\gamma$, $\beta\delta$. Άρα $T(n) = 4 T(n/2) + O(n)$.
Master Theorem: $a = 4,\ b = 2,\ d = 1$, $\log_2 4 = 2 > 1$ $\to$ \textbf{case
3}: $T(n) = O(n^{\log_2 4}) = O(n^2)$.

\textbf{Δεν κερδίσαμε τίποτα!} Ίδια πολυπλοκότητα με τη σχολική μέθοδο.

\subsection{Το έξυπνο τρικ του Karatsuba (1962)}

Παρατηρούμε ότι θέλουμε τρεις ποσότητες: $\alpha\gamma$, $\alpha\delta +
\beta\gamma$ και $\beta\delta$. Αν κάναμε \textbf{τρεις} αντί για τέσσερις
πολλαπλασιασμούς, η αναδρομή θα γινόταν $T(n) = 3T(n/2) + O(n)$ που με Master
Theorem δίνει $O(n^{\log_2 3}) \approx O(n^{1.585})$ --- σημαντική βελτίωση.

\textbf{Το κόλπο:} το \textbf{μεσαίο} $\alpha \delta + \beta \gamma$ μπορεί να
γραφτεί ως διαφορά τριών γινομένων που ήδη ξέρουμε:
\[
(\alpha + \beta)(\gamma + \delta) = \alpha \gamma + \alpha \delta + \beta \gamma
+ \beta \delta
\]
Άρα:
\[
\alpha \delta + \beta \gamma = (\alpha + \beta)(\gamma + \delta) - \alpha \gamma
- \beta \delta
\]
Οι τρεις μόνο πολλαπλασιασμοί που χρειαζόμαστε είναι:
\begin{enumerate}
  \item $z_1 = \alpha \gamma$
  \item $z_2 = \beta \delta$
  \item $z_3 = (\alpha + \beta)(\gamma + \delta)$
\end{enumerate}
Και ο τύπος γίνεται:
\[
x \cdot y = z_1 \cdot 10^n + (z_3 - z_1 - z_2) \cdot 10^{n/2} + z_2
\]

\begin{intuition}
\textbf{Διαισθητικά:} «αγοράζουμε» μια τέταρτη πληροφορία ($\alpha\delta +
\beta\gamma$) χωρίς να κάνουμε τέταρτο πολλαπλασιασμό --- με μερικές προσθέσεις
και αφαιρέσεις που είναι \textbf{πολύ φθηνότερες} ($O(n)$) από τους
πολλαπλασιασμούς. Η ιδέα ανήκει στον Gauss για μιγαδικούς ($(a+bi)(c+di)$) και
την εφάρμοσε στους ακεραίους ο Karatsuba το 1962.
\end{intuition}

\subsection{Ο αλγόριθμος Karatsuba}

\begin{codebox}
\begin{verbatim}
karatsuba(x, y):                  // x, y με n ψηφία
  Εάν n = 1: επίστρεψε x * y       // βάση
  m <- n/2
  Σπάσε x = α*10^m + β
  Σπάσε y = γ*10^m + δ
  z1 <- karatsuba(α, γ)
  z2 <- karatsuba(β, δ)
  z3 <- karatsuba(α + β, γ + δ)
  μέσος <- z3 - z1 - z2
  επίστρεψε z1*10^(2m) + μέσος*10^m + z2
\end{verbatim}
\end{codebox}
\textbf{Σημαντικό:} το $\alpha + \beta$ και $\gamma + \delta$ μπορεί να έχουν
$n/2 + 1$ ψηφία (αν κάνουν κρατούμενο), όχι ακριβώς $n/2$. Αυτή είναι μια
ασήμαντη τεχνική λεπτομέρεια --- δεν επηρεάζει την ασυμπτωτική ανάλυση γιατί
$n/2 + 1 \le n$.

\subsection{Πολυπλοκότητα}

\[
T(n) = 3 T(n/2) + O(n)
\]
Master Theorem: $a = 3,\ b = 2,\ d = 1$. Είναι $\log_2 3 \approx 1.585 > 1$, άρα
$d < \log_b a$ $\to$ \textbf{case 3}:
\[
T(n) = O(n^{\log_2 3}) = O(n^{1.585})
\]

\begin{quoteblock}
\textbf{Θεώρημα (Karatsuba-Ofman, 1962).} Δύο ακέραιοι $n$-ψηφίων μπορούν να
πολλαπλασιαστούν σε $O(n^{1.585})$ βήματα.
\end{quoteblock}

\textbf{Σχόλιο: τι ακολούθησε.} Ο Karatsuba ήταν η \textbf{πρώτη} βελτίωση στο
σχολικό $O(n^2)$ μετά από αιώνες. Άνοιξε ολόκληρο πεδίο έρευνας:

\begin{center}
\begin{tabular}{@{}l l l@{}}
\toprule
Έτος & Αλγόριθμος & Πολυπλοκότητα \\
\midrule
$\sim$1450 BC & Σχολικός                 & $O(n^2)$ \\
1962          & Karatsuba-Ofman          & $O(n^{1.585})$ \\
1963          & Toom-Cook ($k$-way)      & $O(n^{1+1/k})$ \\
1971          & Schönhage-Strassen (FFT) & $O(n \log n \log \log n)$ \\
2007          & Fürer                    & $O(n \log n \cdot 2^{O(\log^* n)})$ \\
2019          & Harvey-van der Hoeven    & $O(n \log n)$ $\checkmark$ \\
\bottomrule
\end{tabular}
\end{center}

Το πρόβλημα \textbf{έκλεισε} το 2019 --- η εργασία Harvey-van der Hoeven έδειξε
ότι ο πολλαπλασιασμός μπορεί να γίνει σε $O(n \log n)$, που έπιανε από καιρό το
θεωρητικό κάτω φράγμα. Στην πράξη, για μικρότερα $n$ το Karatsuba παραμένει το
πιο αποδοτικό λόγω χαμηλών σταθερών --- γι' αυτό χρησιμοποιείται σε βιβλιοθήκες
όπως η GMP.

\begin{recapbox}
\textbf{Τι κρατάμε από αυτή τη διάλεξη}
\begin{enumerate}
  \item \textbf{D\&C ως μέθοδος επίλυσης γενικών προβλημάτων} --- όχι μόνο
  ταξινόμηση. Τρία διαφορετικά προβλήματα, τρεις διαφορετικές «έξυπνες
  συγχωνεύσεις».
  \item \textbf{Μέτρηση αντιστροφών:} η \texttt{merge-and-count} είναι ίδια με
  τη \texttt{merge} της mergesort $+$ μία γραμμή. Αναδρομή ίδια με mergesort
  $\to O(n \log n)$.
  \item \textbf{Κυρίαρχο χρώμα:} η κρίσιμη παρατήρηση («αν είναι κυρίαρχο στην
  αρχική $\to$ είναι κυρίαρχο σε τουλάχιστον μία υποσκακιέρα») περιορίζει τους
  υποψήφιους σε \textbf{4}. Αυτό κάνει τη συγχώνευση γρήγορη $\to O(n^2 \log n)$.
  \item \textbf{Karatsuba:} ο πρώτος αλγόριθμος που έσπασε το $O(n^2)$ για
  πολλαπλασιασμό. \textbf{3} αντί για \textbf{4} αναδρομικές κλήσεις, με μια
  αλγεβρική ταυτότητα $\to O(n^{1.585})$.
  \item \textbf{Σε όλα: το Master Theorem κάνει την ανάλυση μηχανική.} Δες
  $(a, b, d)$, σύγκρινε $d$ με $\log_b a$, διάλεξε case.
\end{enumerate}
\end{recapbox}

\lecturefooter
\end{document}
